\begin{tikzpicture}[
    every label/.append style={anchor=base, yshift=3pt},
    matparens/.style={
      inner sep=3pt,
      left delimiter={(},   % ← metti {[} per quadre, {(} per tonde
      right delimiter={)}   % ← metti {]} per quadre, {)} per tonde
    }
  ]
  % --- PRIMO BLOCCO: ( R  T ) ---
  \node[matparens] (block1) {
    \begin{tikzpicture}
      \node[lipicsGray, matrix={2}{2}, fill=lipicsBulletGray] (R) {};
      %\draw[lipicsGray] (R.north west) -- (R.north east) -- (R.south east) -- cycle;
      \draw[lipicsGray, line width=0.70mm] (R.north west) -- (R.south east) -- cycle;
      %\node[black] at (R.center) {\LARGE$i$};
      \node[font=\LARGE\bfseries, fill=white, draw=gray!30, line width=0.4pt, inner sep=2pt] at (R.center) {$i$};

      \node[gray, matrix={2}{3}, fill=gray, right=2mm of R] (T) {};
      \node[black] at (T.center) {\LARGE$<i$};
    \end{tikzpicture}
  };
  % Freccia
  \node[eqn, right=8mm of block1] (arrow) {$\Rightarrow$};
  % --- SECONDO BLOCCO: ( P  L ) ---
  \node[matparens, right=8mm of arrow] (block2) {
    \begin{tikzpicture}
      \node[lipicsGray, matrix={2}{2}, fill=white] (P) {};
      %\draw[blue, fill=.!10, overlay] (P.north west) -- (P.north east) -- (P.south east) -- cycle;
      %\node[black] at (P.center) [xshift=-2mm] {\LARGE$i$};
      \begin{scope}
        \clip (P.north west) -- (P.south east) -- (P.north east) -- cycle; % triangolo sopra la diagonale
        \fill[lipicsBulletGray] (P.north west) rectangle (P.south east);
      \end{scope}
      \draw[gray, line width=0.70mm] (P.north west) -- (P.south east) -- cycle; % diagonale
      \node[font=\LARGE\bfseries, fill=white, draw=gray!30, line width=0.4pt, inner sep=2pt] at (P.center) [xshift=-2mm] {\LARGE$i$};
      \node[lipicsGray, matrix={2}{2}, fill=none] (P.center) {};

      \node[gray, matrix={2}{3}, fill=gray, right=2mm of P] (L) {};
      \node[black] at (L.center) [xshift=-6mm] {\LARGE$<i$};
    \end{tikzpicture}
  };
\end{tikzpicture}